\documentclass[12pt,a4paper]{article}
\usepackage[margin=2.5cm]{geometry}
\usepackage{amsmath}
\usepackage{hyperref}
\usepackage{parskip}
\usepackage{booktabs}

\title{\textbf{Evaluation}\\[0.3em]
\large Part V(e) --- Modelling Report\\
Asymmetric Sentiment Effects in Earnings Calls}
\author{NLP Course --- Group Project}
\date{}

\begin{document}
\maketitle

\section{Evaluation Protocol}

All models use a strict time-based train/test split to prevent look-ahead bias:

\begin{itemize}
  \item \textbf{Training}: years 2016--2018 ($n \approx 155$ events)
  \item \textbf{Test}: years 2019--2020 ($n \approx 187$ events)
\end{itemize}

No test-period data enters feature vocabulary construction, scaling, or
hyperparameter selection.  This is essential because many events in the test window
occurred during or after COVID-19 market disruptions, and any data leakage would
make the evaluation unreliable.

\section{Regression Metrics}

\subsection{OOS $R^2$}

Following Campbell and Thompson (2008), out-of-sample $R^2$ is:
\[
  \text{OOS }R^2 = 1 - \frac{\sum_{t \in \text{test}} (y_t - \hat{y}_t)^2}
                              {\sum_{t \in \text{test}} (y_t - \bar{y}_{\text{train}})^2}
\]
A positive value means the model beats the naive historical-mean forecast.
This is a stricter benchmark than in-sample $R^2$, which can be inflated by overfitting.

\subsection{MAE}

Mean Absolute Error in CAR units gives an intuitive sense of typical prediction error.

\subsection{Wald Test p-value}

For OLS models, we test $H_0: \beta_{\text{neg}} + \beta_{\text{pos}} = 0$ using a
Wald test with HC3 standard errors.  Rejection at $p < 0.10$ supports the asymmetry
hypothesis.

\section{Regression Results}

\begin{center}
\begin{tabular}{lr}
\toprule
Model & OOS $R^2$ \\
\midrule
Null (historical mean)      & $0.000$ \\
\textbf{LM Lexicon [OLS]}   & $\mathbf{+0.057}$ \\
Word2Vec [Ridge]             & $+0.050$ \\
TF-IDF [RF]                  & $+0.045$ \\
TF-IDF [Ridge]               & $+0.043$ \\
LM + Controls [Ridge]        & $-0.044$ \\
LM + Controls [OLS]          & $-0.086$ \\
TF-IDF [OLS]                 & $-0.469$ \\
\bottomrule
\end{tabular}
\end{center}

The LM lexicon achieves the best OOS $R^2$ of $+0.057$.  Models with year dummy
controls overfit with only $\sim$155 training observations, even with Ridge regularization.
The Wald test rejects $H_0$ at $p = 0.042$ for LM + Controls [OLS],
with $\hat\beta_{\text{neg}} = -8.92$ vs.\ $\hat\beta_{\text{pos}} = +3.31$.

\section{Classification Metrics}

Accuracy, macro F1, and AUC are reported on the 2019--2020 test set.
F1 and AUC are more informative than accuracy because the class distribution
in the test window skews toward positive CARs (bull market conditions in 2019--2020).

\begin{center}
\begin{tabular}{lrrr}
\toprule
Model & Accuracy & F1 & AUC \\
\midrule
TF-IDF + LM + Controls [RF] & $59.6\%$ & $0.553$ & $\mathbf{0.600}$ \\
TF-IDF [NaiveBayes]         & $\mathbf{59.6\%}$ & $0.632$ & $0.588$ \\
TF-IDF [RF]                  & $51.9\%$ & $0.468$ & $0.563$ \\
TF-IDF [LogReg]              & $51.9\%$ & $0.590$ & $0.578$ \\
LM Lexicon [LogReg]          & $45.5\%$ & $0.625$ & $0.508$ \\
\bottomrule
\end{tabular}
\end{center}

Best AUC is $0.600$; best accuracy is $59.6\%$.
All TF-IDF models beat the 50\% random baseline in accuracy.

\section{Cross-Validation}

An expanding-window CV (train on 2016$\rightarrow$Y, test on Y+1 for Y=2017--2019)
confirms that LM lexicon features maintain positive pooled OOS $R^2$ across folds,
while TF-IDF and control-heavy models remain more variable across test years.

\section{Code Reference}

\begin{itemize}
  \item \texttt{src/nasdaq\_nlp/evaluation/} --- OOS $R^2$, Wald test, AUC utilities.
  \item \texttt{notebooks/03\_modeling.ipynb} --- full evaluation runs.
  \item \texttt{notebooks/04\_results.ipynb} --- result visualization and interpretation.
\end{itemize}

\end{document}
